% ============================================================
%  Shared preamble: USABO, IBO and Beyond -- Biology Compendium
% ============================================================
\documentclass[10pt,letterpaper]{article}

\usepackage[T1]{fontenc}
\usepackage[utf8]{inputenc}
\usepackage[letterpaper,left=1.10in,right=1.10in,top=1.00in,bottom=0.95in,
            headsep=17pt,footskip=28pt]{geometry}

% ---- Fonts -------------------------------------------------
\usepackage{mathpazo}
\usepackage[scaled=0.90]{helvet}
\usepackage{courier}
\usepackage[expansion=false]{microtype}
\linespread{1.22}
\setlength{\parskip}{8.5pt}
\setlength{\parindent}{0pt}
\setlength{\emergencystretch}{3em}

% ---- Mathematics and chemistry -----------------------------
\usepackage{amsmath,amssymb,amsthm,mathtools}
\usepackage{bm}
\usepackage[version=4]{mhchem}
\allowdisplaybreaks

% ---- Graphics and tables -----------------------------------
\usepackage{graphicx}
\usepackage{booktabs}
\usepackage{array}
\usepackage{multicol}
\usepackage{adjustbox}
\usepackage{longtable}
\usepackage{enumitem}
\setlist{nosep,topsep=3pt,itemsep=3pt,parsep=0pt,leftmargin=15pt}

\usepackage{tikz}
\usetikzlibrary{calc,intersections,angles,quotes,patterns,arrows.meta,positioning,decorations.pathreplacing}

% ---- Colours -----------------------------------------------
\usepackage{xcolor}
\definecolor{bioink}{HTML}{16332B}
\definecolor{biogreen}{HTML}{2C6E54}
\definecolor{bioline}{HTML}{C5D2CB}
\definecolor{biogrey}{HTML}{F3F6F4}
\definecolor{biogold}{HTML}{9C6B1E}
\definecolor{bioteal}{HTML}{1F6F6B}
\definecolor{biored}{HTML}{9B2C2C}
\definecolor{bioplum}{HTML}{6B3A63}
\definecolor{bionum}{HTML}{9BA8A1}
\definecolor{tiercore}{HTML}{2C6E54}
\definecolor{tierbeyond}{HTML}{6B5B3E}
\definecolor{tiercat}{HTML}{3A5A7A}

% ---- Boxes -------------------------------------------------
\usepackage[most]{tcolorbox}
\tcbset{biobase/.style={enhanced,arc=0pt,outer arc=0pt,
  left=11pt,right=0pt,top=5pt,bottom=5pt,
  boxrule=0pt,colframe=white,colback=white,
  fonttitle=\sffamily\bfseries\footnotesize,
  toptitle=0pt,bottomtitle=3pt,
  before skip=13pt,after skip=13pt}}

\newtcolorbox{keybox}[1][]{biobase,breakable,colback=biogrey,colbacktitle=biogrey,
  coltitle=biogreen,borderline west={2.2pt}{0pt}{biogreen},#1}
\newtcolorbox{thmbox}[1][]{biobase,breakable,colback=white,colbacktitle=white,
  coltitle=bioink,borderline west={2.2pt}{0pt}{bioink},#1}
\newtcolorbox{defbox}[1][]{biobase,breakable,colback=white,colbacktitle=white,
  coltitle=bioteal,borderline west={2.2pt}{0pt}{bioteal},#1}
\newtcolorbox{warnbox}[1][]{biobase,breakable,colback=white,colbacktitle=white,
  coltitle=biored,borderline west={2.2pt}{0pt}{biored},#1}
\newtcolorbox{tipbox}[1][]{biobase,breakable,colback=white,colbacktitle=white,
  coltitle=biogreen,borderline west={2.2pt}{0pt}{biogreen},#1}
\newtcolorbox{obsbox}[1][]{biobase,breakable,colback=white,colbacktitle=white,
  coltitle=biogold,borderline west={2.2pt}{0pt}{biogold},#1}
\newtcolorbox{calcbox}[1][]{biobase,breakable,colback=biogrey,colbacktitle=biogrey,
  coltitle=bioink,borderline west={2.2pt}{0pt}{bionum},#1}
\newtcolorbox{expbox}[1][]{biobase,breakable,colback=white,colbacktitle=white,
  coltitle=bioplum,borderline west={2.2pt}{0pt}{bioplum},#1}

% Non-breakable clones, for use inside minipages and tables
\newtcolorbox{keyboxnb}[1][]{biobase,colback=biogrey,colbacktitle=biogrey,
  coltitle=biogreen,borderline west={2.2pt}{0pt}{biogreen},#1}
\newtcolorbox{thmboxnb}[1][]{biobase,colback=white,colbacktitle=white,
  coltitle=bioink,borderline west={2.2pt}{0pt}{bioink},#1}
\newtcolorbox{calcboxnb}[1][]{biobase,colback=biogrey,colbacktitle=biogrey,
  coltitle=bioink,borderline west={2.2pt}{0pt}{bionum},#1}

% ---- Index -------------------------------------------------
\usepackage{imakeidx}

% ---- Sectioning --------------------------------------------
\usepackage{titlesec}
\usepackage{titletoc}
\usepackage[hidelinks]{hyperref}
\hypersetup{colorlinks=true,linkcolor=biogreen,urlcolor=biogreen,citecolor=biogreen,
  bookmarksnumbered=true}

\titleformat{\section}{\sffamily\bfseries\large\color{bioink}}{\thesection}{0.6em}{}
\titleformat{\subsection}{\sffamily\bfseries\normalsize\color{biogreen}}{\thesubsection}{0.6em}{}
\titlespacing*{\section}{0pt}{16pt}{6pt}
\titlespacing*{\subsection}{0pt}{12pt}{4pt}

% ---- Running heads -----------------------------------------
\usepackage{fancyhdr}
\pagestyle{fancy}
\fancyhf{}
\renewcommand{\headrulewidth}{0pt}
\fancyhead[L]{\sffamily\footnotesize\color{bionum}\leftmark}
\fancyfoot[C]{\sffamily\footnotesize\color{bionum}\thepage}

% ---- Small utilities ---------------------------------------
\newcommand{\tm}[1]{\textbf{#1}}
\newcommand{\dg}{^{\circ}}
\newcommand{\degc}{\,^{\circ}\mathrm{C}}
\newcommand{\sn}[1]{\textit{#1}}
\newcommand{\ttitle}[1]{\textbf{\textcolor{bioink}{#1}}}
\newcommand{\secchip}[1]{%
  \noindent{\sffamily\bfseries\fontsize{13.5}{16.5}\selectfont\color{bioink}\raggedright #1\par}%
  \vspace{4pt}%
  \noindent{\color{biogreen}\rule{\textwidth}{1.1pt}}\par}

% Named result: index entry plus a label the catalogue can point at
\newcommand{\nr}[2]{\index{#2}\label{nr:#1}}
\newcommand{\nrrow}[2]{#2\dotfill\pageref{nr:#1}\\}

% ---- Parts -------------------------------------------------
\newcounter{biopart}
\newcommand{\prepart}{}
\newcommand{\pref}[1]{Part~\ref{part:#1}}
\newcommand{\prefn}[1]{\ref{part:#1}}

\newcommand{\parttitle}[3][]{%
  \clearpage
  \refstepcounter{biopart}%
  \ifx\relax#1\relax\else\label{part:#1}\fi
  \markboth{#2}{#2}%
  \thispagestyle{fancy}%
  \vspace*{6pt}%
  \noindent{\color{bioline}\rule{\textwidth}{1pt}}\\[6pt]
  \noindent{\sffamily\bfseries\footnotesize\color{bionum}PART \thebiopart}\\[8pt]
  \noindent\begin{minipage}{\textwidth}\raggedright
    {\sffamily\bfseries\fontsize{22}{26}\selectfont\color{bioink}%
     \hyphenchar\font=-1\hyphenpenalty=10000\exhyphenpenalty=10000 #2\par}%
  \end{minipage}\\[10pt]
  \noindent{\color{biogreen}\rule{\textwidth}{1.4pt}}\\[8pt]
  \noindent\begin{minipage}{\textwidth}\raggedright
    {\fontsize{10.5}{14}\selectfont\color{biogreen} #3\par}%
  \end{minipage}
  \vspace{10pt}\par
  \addcontentsline{toc}{section}{\protect\numberline{\thebiopart}#2}%
}

% ---- Tier dividers -----------------------------------------
\newcommand{\tierpartline}[1]{\noindent\hangindent=12pt {\color{biogreen}\textbullet}\ #1\par\vspace{1pt}}
\newcommand{\tierdivider}[4]{%
  \clearpage
  \thispagestyle{empty}%
  \vspace*{0.08\textheight}
  \noindent{\color{tier#1}\rule{\textwidth}{5pt}}\\[3pt]
  \noindent{\color{bioline}\rule{\textwidth}{1pt}}\\[26pt]
  \noindent\begin{minipage}{\textwidth}\raggedright
    {\sffamily\bfseries\fontsize{30}{36}\selectfont\color{bioink} #2\par}%
  \end{minipage}\\[16pt]
  \noindent{\color{bioline}\rule{\textwidth}{1pt}}\\[12pt]
  \noindent\begin{minipage}{\textwidth}\raggedright
    {\fontsize{11.5}{16}\selectfont\color{biogreen} #3\par}%
  \end{minipage}\\[20pt]
  {\footnotesize\begin{multicols}{2}#4\end{multicols}}
  \vfill
  \clearpage
}

% ---- Table safety: never let a tabular run past the measure --
\usepackage{etoolbox}
\BeforeBeginEnvironment{tabular}{\begin{adjustbox}{max width=\linewidth}}
\AfterEndEnvironment{tabular}{\end{adjustbox}}

% Part display names, so a cross reference into a part that is not in this volume degrades
% to the part's name in italics instead of a dangling question mark.
\makeatletter
\expandafter\def\csname pname@foundations\endcsname{Foundations, Scale, and the Language of Biology}
\expandafter\def\csname pname@chemistry\endcsname{The Chemistry of Life}
\expandafter\def\csname pname@macromolecules\endcsname{Biological Macromolecules}
\expandafter\def\csname pname@enzymes\endcsname{Enzymes, Energetics, and Metabolic Control}
\expandafter\def\csname pname@cellstructure\endcsname{Cell Structure and the Organelles}
\expandafter\def\csname pname@membranes\endcsname{Membranes and Membrane Transport}
\expandafter\def\csname pname@respiration\endcsname{Cellular Respiration and Fermentation}
\expandafter\def\csname pname@photosynthesis\endcsname{Photosynthesis}
\expandafter\def\csname pname@signaling\endcsname{Cell Signaling and Signal Transduction}
\expandafter\def\csname pname@cytoskeleton\endcsname{The Cytoskeleton, Motility, and the Extracellular Matrix}
\expandafter\def\csname pname@cellcycle\endcsname{The Cell Cycle, Mitosis, and Meiosis}
\expandafter\def\csname pname@cancer\endcsname{Cancer and the Control of Cell Number}
\expandafter\def\csname pname@replication\endcsname{DNA Replication, Repair, and Chromosome Structure}
\expandafter\def\csname pname@transcription\endcsname{Transcription, RNA Processing, and the Genetic Code}
\expandafter\def\csname pname@translation\endcsname{Translation and Protein Targeting}
\expandafter\def\csname pname@generegulation\endcsname{Regulation of Gene Expression and Epigenetics}
\expandafter\def\csname pname@biotech\endcsname{Biotechnology and Molecular Laboratory Methods}
\expandafter\def\csname pname@mendelian\endcsname{Mendelian Genetics, Pedigrees, and Probability}
\expandafter\def\csname pname@linkage\endcsname{Linkage, Mapping, and Chromosomal Genetics}
\expandafter\def\csname pname@popgen\endcsname{Population Genetics and Hardy-Weinberg}
\expandafter\def\csname pname@evolution\endcsname{Evolutionary Mechanisms and Natural Selection}
\expandafter\def\csname pname@macroevolution\endcsname{Speciation, Macroevolution, and the History of Life}
\expandafter\def\csname pname@phylogenetics\endcsname{Phylogenetics and Biosystematics}
\expandafter\def\csname pname@prokaryotes\endcsname{Bacteria, Archaea, and the Prokaryotic World}
\expandafter\def\csname pname@viruses\endcsname{Viruses and Subcellular Agents}
\expandafter\def\csname pname@protists\endcsname{Protists and the Origin of Eukaryotes}
\expandafter\def\csname pname@fungi\endcsname{Fungi}
\expandafter\def\csname pname@plantdiversity\endcsname{Plant Diversity, Life Cycles, and Reproduction}
\expandafter\def\csname pname@plantanatomy\endcsname{Plant Anatomy, Tissues, and Growth}
\expandafter\def\csname pname@planttransport\endcsname{Plant Transport, Nutrition, and Water Relations}
\expandafter\def\csname pname@planthormones\endcsname{Plant Hormones, Signals, and Development}
\expandafter\def\csname pname@animaldiversity\endcsname{Animal Diversity and Body Plans}
\expandafter\def\csname pname@musculoskeletal\endcsname{Animal Tissues, Skeleton, Muscle, and Movement}
\expandafter\def\csname pname@digestion\endcsname{Animal Nutrition and Digestion}
\expandafter\def\csname pname@gasexchange\endcsname{Gas Exchange and Respiratory Systems}
\expandafter\def\csname pname@circulation\endcsname{Circulation, Blood, and the Heart}
\expandafter\def\csname pname@immunology\endcsname{Immunology}
\expandafter\def\csname pname@osmoregulation\endcsname{Osmoregulation, Excretion, and Thermoregulation}
\expandafter\def\csname pname@nervous\endcsname{The Nervous System and Neurophysiology}
\expandafter\def\csname pname@sensory\endcsname{Sensory Biology}
\expandafter\def\csname pname@endocrine\endcsname{The Endocrine System}
\expandafter\def\csname pname@development\endcsname{Animal Reproduction and Development}
\expandafter\def\csname pname@ethology\endcsname{Ethology and Animal Behavior}
\expandafter\def\csname pname@popecology\endcsname{Population and Community Ecology}
\expandafter\def\csname pname@ecosystems\endcsname{Ecosystems, Biomes, Biogeochemistry, and Conservation}
\expandafter\def\csname pname@examcraft\endcsname{Exam Craft: The Open Exam, the Semifinal, and the IBO}
\expandafter\def\csname pname@devbio\endcsname{Developmental Biology in Depth}
\expandafter\def\csname pname@neuro\endcsname{Neurobiology in Depth}
\expandafter\def\csname pname@immuno\endcsname{Immunology in Depth}
\expandafter\def\csname pname@microbio\endcsname{Microbiology and Virology in Depth}
\expandafter\def\csname pname@compphys\endcsname{Comparative and Clinical Physiology}
\expandafter\def\csname pname@structural\endcsname{Structural Biology and Protein Science}
\expandafter\def\csname pname@systems\endcsname{Systems Biology and Gene Networks}
\expandafter\def\csname pname@techniques\endcsname{Molecular Techniques in Depth}
\expandafter\def\csname pname@evogenomics\endcsname{Evolutionary Genomics and Molecular Evolution}
\expandafter\def\csname pname@practical\endcsname{The Practical Exam: Lab Skills, Histology, and Data Interpretation}
\expandafter\def\csname pname@f-units\endcsname{Units, Scaling, Dimensional Analysis, and Magnitudes}
\expandafter\def\csname pname@f-biostats\endcsname{Biostatistics: Distributions, Estimation, and Error}
\expandafter\def\csname pname@f-tests\endcsname{Hypothesis Tests, Chi-Square, and Goodness of Fit}
\expandafter\def\csname pname@f-probability\endcsname{Probability, Pedigrees, and Bayesian Genetics}
\expandafter\def\csname pname@f-mendelian\endcsname{Mendelian and Linkage Calculations}
\expandafter\def\csname pname@f-hardyweinberg\endcsname{Hardy-Weinberg and Population Genetics Equations}
\expandafter\def\csname pname@f-selection\endcsname{Selection, Drift, Mutation, Migration, and Inbreeding}
\expandafter\def\csname pname@f-kinetics\endcsname{Enzyme Kinetics Equations}
\expandafter\def\csname pname@f-thermo\endcsname{Thermodynamics, Bioenergetics, and Redox Calculations}
\expandafter\def\csname pname@f-membrane\endcsname{Membrane Potentials: Nernst, Goldman, and Cable Theory}
\expandafter\def\csname pname@f-transport\endcsname{Diffusion, Osmosis, Flow, and Water Potential}
\expandafter\def\csname pname@f-cardioresp\endcsname{Respiratory and Cardiovascular Calculations}
\expandafter\def\csname pname@f-renal\endcsname{Renal, Osmoregulatory, and Acid-Base Calculations}
\expandafter\def\csname pname@f-metabolic\endcsname{Metabolic Rate, Allometry, and Thermal Physiology}
\expandafter\def\csname pname@f-plantcalc\endcsname{Photosynthesis, Plant Water, and Growth Calculations}
\expandafter\def\csname pname@f-popgrowth\endcsname{Population Growth, Life Tables, and Demography}
\expandafter\def\csname pname@f-ecosystem\endcsname{Community, Diversity, and Ecosystem Formulas}
\expandafter\def\csname pname@f-epidemiology\endcsname{Epidemiology and Disease Dynamics}
\expandafter\def\csname pname@f-labcalc\endcsname{Laboratory, Molecular, and Bioinformatic Calculations}
\expandafter\def\csname pname@f-calculus\endcsname{Calculus, Differential Equations, and Modelling Shortcuts}
\expandafter\def\csname pname@f-allformulas\endcsname{Every Formula, Constant, and Normal Value}
\expandafter\def\csname pname@c-laws\endcsname{Named Laws, Rules, and Principles of Biology}
\expandafter\def\csname pname@c-pathways\endcsname{Metabolic Pathways and Cycles}
\expandafter\def\csname pname@c-molecules\endcsname{Molecules, Enzymes, Cofactors, and Vitamins}
\expandafter\def\csname pname@c-proteins\endcsname{Proteins, Receptors, Channels, and Transporters}
\expandafter\def\csname pname@c-genes\endcsname{Genes, Mutations, and Genetic Disorders}
\expandafter\def\csname pname@c-hormones\endcsname{Hormones and Signaling Molecules}
\expandafter\def\csname pname@c-experiments\endcsname{Model Organisms and Landmark Experiments}
\expandafter\def\csname pname@c-taxonomy\endcsname{Taxonomy: Every Phylum and Major Clade}
\expandafter\def\csname pname@c-anatomy\endcsname{Anatomy, Histology, and Normal Values}
\expandafter\def\csname pname@c-terminology\endcsname{Terminology, Roots, and the Timeline of Biology}
\newcommand{\partnameof}[1]{%
  \ifcsname pname@#1\endcsname\csname pname@#1\endcsname\else this book\fi}
\renewcommand{\pref}[1]{\@ifundefined{r@part:#1}{\emph{\partnameof{#1}}}{Part~\ref{part:#1}}}
\renewcommand{\prefn}[1]{\@ifundefined{r@part:#1}{\emph{\partnameof{#1}}}{\ref{part:#1}}}
\renewcommand{\nrrow}[2]{\@ifundefined{r@nr:#1}{#2}{#2\dotfill\pageref{nr:#1}}\\}
\makeatother

\begin{document}
\begin{titlepage}\thispagestyle{empty}
\vspace*{0.14\textheight}
\noindent{\color{tiercat}\rule{\textwidth}{5pt}}\\[3pt]
\noindent{\color{bioline}\rule{\textwidth}{1pt}}\\[26pt]
\noindent{\sffamily\bfseries\fontsize{11}{14}\selectfont\color{bionum}START HERE}\\[12pt]
\noindent\begin{minipage}{\textwidth}\raggedright
{\sffamily\bfseries\fontsize{30}{34}\selectfont\color{bioink}%
 \hyphenchar\font=-1 How to Use These Formulas\par}\end{minipage}\\[16pt]
\noindent{\color{biogreen}\rule{\textwidth}{1.4pt}}\\[10pt]
\noindent\begin{minipage}{\textwidth}\raggedright
{\fontsize{10.5}{14}\selectfont\color{biogreen}%
Which equation a question is asking for, how to recognise it from the wording, what to
memorise against what to rebuild, and the checks that catch a wrong answer before you
write it down.\par}\end{minipage}
\vfill
\noindent{\color{bioline}\rule{\textwidth}{1pt}}\\[3pt]
\noindent{\color{tiercat}\rule{\textwidth}{5pt}}
\end{titlepage}
\tableofcontents
\clearpage

\section{What is in this zip}

\begin{keybox}[title={The two kinds of file}]
\begin{itemize}
\item \tm{The index of every formula}, \texttt{00-Every-Formula-Indexed.pdf}. Every
      equation, law, constant and normal value in the compendium, stated once with its
      symbols defined and nothing derived, ordered by subject with an alphabetical
      index at the back. Revise from this file and look things up in it.
\item \tm{The formula volumes}, in \texttt{01-Volumes}. One per quantitative area. Part
      one derives every formula in the area, states the conditions under which it holds,
      and works the calculations with the arithmetic shown. Part two is the bare formula
      sheet for the same material.
\end{itemize}
The companion zip, \texttt{USABO-Biology-Concepts}, holds the biology these equations
describe. Work the two together: a Nernst potential means nothing without the ion
gradients, and the ion gradients are hard to remember without the calculation.
\end{keybox}

\section{Why a biology exam has formulas at all}

Every stage of this competition asks for arithmetic, and it is the part most candidates
have not practised. The Open Exam hides calculations inside multiple choice options, so
a wrong unit conversion looks exactly like a plausible distractor. The Semifinal Exam
asks for them openly, in data analysis questions where the numbers come off a figure.
The IBO practicals are almost entirely calculation: a magnification, a dilution, a rate,
a statistic, an error estimate.

\begin{obsbox}[title={WORTH NOTICING: the recurring seven}]
Across past papers, seven calculations recur more than any others: magnification and
true size from a micrograph; a dilution or a solution preparation; Hardy-Weinberg allele
frequencies; a chi-square goodness of fit; a rate from a graph's gradient; a Nernst or
Goldman potential; and a population growth or mark-recapture estimate. If your practice
time is short, those seven are where it goes.
\end{obsbox}

\section{Recognising which formula is wanted}

\begin{keybox}[title={The wording-to-equation key}]
\begin{itemize}
\item ``What fraction of the population carries'' $\Rightarrow$ Hardy-Weinberg, and
      usually $2pq$ rather than $q^2$.
\item ``Do the observed numbers differ from expected'' $\Rightarrow$ chi-square, and
      state the degrees of freedom before computing.
\item ``Equilibrium potential'' or ``at what voltage does the ion stop moving''
      $\Rightarrow$ Nernst. ``Resting potential'' with several permeabilities given
      $\Rightarrow$ Goldman.
\item ``Per unit time, across a surface, given a gradient'' $\Rightarrow$ Fick.
\item ``Half maximal'' anything $\Rightarrow$ $K_m$, or $P_{50}$, or an $EC_{50}$: the
      same algebra each time.
\item ``How much of the plasma is cleared'' $\Rightarrow$ clearance, $C=UV/P$.
\item ``Water moves which way'' $\Rightarrow$ water potential, and check the sign twice.
\item ``Rate doubles every ten degrees'' $\Rightarrow$ $Q_{10}$.
\item ``Estimate the population'' from two captures $\Rightarrow$ Lincoln index.
\item ``Energy available to the next trophic level'' $\Rightarrow$ transfer efficiency,
      and do not assume exactly ten per cent unless told to.
\end{itemize}
\end{keybox}

\section{What to memorise and what to rebuild}

\begin{keybox}[title={Memorise these outright}]
Hardy-Weinberg in both forms; the chi-square statistic; Nernst at body temperature;
Fick's law; the Michaelis-Menten equation and its double reciprocal form;
Henderson-Hasselbalch; water potential and its components; clearance; cardiac output;
the exponential and logistic growth equations; the Lincoln index; magnification and true
size; the Beer-Lambert law; $Q_{10}$; and the gas constant, Faraday constant and
Avogadro number.
\end{keybox}

\begin{keybox}[title={Rebuild these from a principle, never memorise}]
Anything with several variants: the inhibition equations, which all follow from where
the inhibitor binds; the acid-base compensations, which follow from
Henderson-Hasselbalch; the selection recursions, which follow from a fitness table; the
diversity indices, which follow from what each is trying to weight; and every
differential equation model, which follows from writing down gains minus losses.
\end{keybox}

\section{The checks that catch a wrong answer}

\begin{tipbox}[title={TECHNIQUE: four audits, in this order}]
\begin{enumerate}
\item \tm{Units.} Carry them through every line. A clearance must come out in
      $\mathrm{mL\,min^{-1}}$; if it does not, the error is upstream of the arithmetic.
\item \tm{Direction and sign.} Does water move the way your answer says? Is the
      potential negative for an anion? Is the free energy negative for a spontaneous
      step? Most lost marks are sign errors, not algebra.
\item \tm{Magnitude.} Compare with a value you know. A resting potential near
      $-70\ \mathrm{mV}$, a GFR near $125\ \mathrm{mL\,min^{-1}}$, a bacterial doubling
      time of twenty minutes. An answer three orders of magnitude out is a unit prefix
      slip.
\item \tm{Limit case.} Set a variable to zero or to infinity and see whether the formula
      still says something sensible. This catches an inverted fraction instantly.
\end{enumerate}
\end{tipbox}

\begin{warnbox}[title={TRAP: the six errors that cost the most marks}]
\begin{itemize}
\item Using $q^2$ where the question asked for carriers, which is $2pq$.
\item Forgetting that Hardy-Weinberg for an X linked locus has different male and
      female expectations.
\item Mixing $\mathrm{mM}$ with $\mathrm{M}$ inside a log, which shifts the answer by
      exactly a factor of $\ln 1000$.
\item Reading a semilog axis as if it were linear.
\item Quoting standard error as standard deviation, or the reverse, when a figure's
      error bars are unlabelled. Say which you are assuming.
\item Rounding at every intermediate step instead of the last one.
\end{itemize}
\end{warnbox}

\section{An order to work through the formula volumes}

\begin{keybox}[title={Dependency order}]
\begin{enumerate}
\item Units, scaling and dimensional analysis. Everything else assumes it.
\item Biostatistics, then hypothesis tests and chi-square. These are used by every
      later volume and by every data question in the exam.
\item Probability and pedigrees, then Mendelian and linkage calculations, then
      Hardy-Weinberg, then selection and drift, in that order. Each rests on the last.
\item Thermodynamics and redox, then enzyme kinetics. Kinetics needs the free energy
      language.
\item Diffusion and water potential, then membrane potentials. Both are Fick's law
      wearing different hats.
\item Cardiorespiratory, then renal and acid-base, then metabolic rate and thermal
      physiology. The physiology block, in the order the concepts zip teaches it.
\item Photosynthesis and plant calculations.
\item Population growth, then community and ecosystem formulas, then epidemiology. The
      last is the same mathematics as the first with a different label.
\item Laboratory and bioinformatic calculations, at any point, in short sittings. These
      are drills rather than theory.
\item Calculus and modelling last. It explains where the earlier models came from, and
      it is the volume that makes the rest feel like one subject instead of twenty.
\end{enumerate}
\end{keybox}

\begin{tipbox}[title={How to practise}]
Do not read a worked calculation. Cover the solution, do it, then compare. A formula
volume read passively transfers almost nothing, because the skill being tested is
execution under time pressure, and execution is only trained by executing. Twenty
minutes of covered-solution practice beats two hours of reading, every time.
\end{tipbox}
\end{document}
